\documentclass[10pt,a4paper]{amsart}
\usepackage[margin=19mm]{geometry}
\usepackage{amsmath,amssymb,mathtools}
\usepackage[scr=rsfs,cal=euler]{mathalfa}
\usepackage{xfrac}
\usepackage[unicode,hidelinks,bookmarks=false]{hyperref}
\newcommand{\ie}{i.e.\ }
\newcommand{\cf}{cf.\ }
\newcommand{\resp}{respectively\ }
\newcommand{\Subsection}{\subsection}
\providecommand{\nobreakdash}{\nobreak\hbox{-}}
\setlength{\emergencystretch}{2em}
\multlinegap0pt
\usepackage{amssymb}
\usepackage{mathtools}
\usepackage{tikz}
\newtheorem{proposition}{Proposition}
\title{On non-positive Weyl connections on Lie groups}
\author{Tomasz Dudek, Aleksy Tralle}
\date{Source: arXiv:2607.17210v1 (CC BY 4.0)}
\begin{document}
\maketitle

\noindent\emph{Source and licence.} On non-positive Weyl connections on Lie groups. Version arXiv:2607.17210v1 (CC BY 4.0), distributed by arXiv under the Creative Commons Attribution 4.0 International licence, http://creativecommons.org/licenses/by/4.0/, as recorded in the arXiv licence metadata for this exact version and shown on the abstract page https://arxiv.org/abs/2607.17210v1. Full licence notice: \texttt{licenses/arxiv-2607.17210-CC-BY-4.0.txt}.\par\smallskip

\noindent\textit{Editorial scope.} Counted unit: the article's proposition prop:hat-k (source0.tex lines 127--132). The article's complete original proof of it is delivered with it (source0.tex lines 133--138, one \texttt{proof} environment of 6 lines). No mathematics is written, repaired or re-derived: the statement and the proof are the article's own lines, in the article's own notation, and no retained line is substituted or re-flowed.  No further source-local context is needed: every cross-reference of every retained line resolves inside the two delivered spans.  Nothing was omitted from any retained span, so the package carries no omission record.  The retained lines cite 1 works, whose entries are transcribed verbatim from the article's own bibliography source in the e-print and delivered with short editorial labels recorded in the item's reference mapping.  No retained line contains a figure environment, an image-inclusion command or drawing code, so no image is delivered and the package declares no asset dependency.  Editorial formatting only, in addition: an AMS \texttt{amsart} preamble that restates the article's own declarations and macros used by the retained lines, namely the environments \texttt{proposition} on separate counters, together with the standard packages those lines need.  The retained environments print in this document as Proposition 1 in order of appearance, whereas the article prints the same items with the numbering of its own full document.  The complete native source of the article as distributed in the arXiv e-print for this exact version is bundled unchanged as \texttt{sources/205547/source0.tex}.
\begin{proposition}\label{prop:hat-k} Let $(G,\langle -,-\rangle)$ be a Lie group endowed with a left-invariant Weyl connection determined by $E\in\mathfrak{g}$. The Weyl curvature of this connection is given by the formula
 \begin{equation}\widehat{K}(X,Y) =$$
$$ K(X,Y)  + \langle {Y},{E}\rangle^2 + \langle {X},{E}\rangle^2 - \langle{E},{E}\rangle - \langle{Y},{\alpha(Y,E)}\rangle  - \langle{\alpha(X,E)},{X}\rangle,
\end{equation}
for any unit orthogonal $X,Y$.
\end{proposition}

\begin{proof} Use the following formula for the Weyl curvature \cite{W} (Proposition 3.2):
$$\widehat{R}_a(X,Y)Z=R(X,Y)Z+\langle Z,E\rangle(\langle Y,E\rangle X-\langle X,E\rangle Y)$$
$$+(\langle Z,Y\rangle \langle X,E\rangle-\langle Z,X\rangle\langle Y,E\rangle)E+E^2(\langle Z,X\rangle Y-\langle Z,Y\rangle X)$$
$$+\langle Z,\nabla_XE\rangle Y-\langle Z,\nabla_YE\rangle X+\langle Z,X\rangle\nabla_YE-\langle Z,Y\rangle\nabla_XE$$ 
  together with assumptions $Z=Y$, $\langle X,X\rangle=\langle Y,Y\rangle=1$, $\langle X,Y\rangle=0$ and the formula $(\nabla_UV)_o=\alpha(U,V)$.
\end{proof}

\begin{thebibliography}{99}
\bibitem[W]{W} M. P. Wojtkowski, {\it W-flows on Weyl manifolds and Gaussian theromostats.} J. Math. Pure Appl., 10(2000),953-974
\end{thebibliography}
\end{document}
